\chapter{Material adicional}
\label{ap:material}

\section{Instalación del entorno}
Instrucciones para reproducir el entorno de desarrollo:
\begin{lstlisting}[language=bash, numbers=none]
python -m venv .venv
source .venv/bin/activate        # En Windows: .venv\Scripts\activate
pip install -r requirements.txt
\end{lstlisting}

\section{Descripción de los campos de CBIS-DDSM}

\subsection{Manifiesto de TCIA (incluido en el repositorio)}
El fichero \texttt{CBIS-DDSM-All-doiJNLP-zzWs5zfZ-nbia-digest.xlsx} es el
manifiesto que genera TCIA al exportar una colección: identifica las series
DICOM disponibles pero \textbf{no incluye la etiqueta de patología}. Sus
campos más relevantes para el \textit{pipeline} son:

\begin{table}[H]
\centering
\caption{Campos relevantes del manifiesto de TCIA.}
\label{tab:campos-manifiesto}
\begin{tabular}{@{}p{3.2cm}p{11cm}@{}}
\toprule
\textbf{Campo} & \textbf{Descripción} \\
\midrule
\texttt{Patient ID} & Identificador largo, p.\,ej.
  \texttt{Mass-Training\_P\_01009\_RIGHT\_CC\_1}. \texttt{src.data.load\_manifest}
  extrae de aquí el id real de paciente (\texttt{P\_01009}) con una expresión
  regular, base de la partición por paciente. \\
\texttt{Series Description} & Tipo de imagen de la serie: \texttt{full mammogram
  images}, \texttt{ROI mask images} o \texttt{cropped images}. \\
\texttt{Modality} & Modalidad DICOM; \texttt{MG} (mamografía) en todos los casos. \\
\texttt{Body Part Examined} & Región anatómica (\texttt{BREAST}). \\
\texttt{Image Count} & Número de imágenes DICOM de la serie. \\
\texttt{File Size} & Tamaño en bytes de la serie (para estimar la descarga total). \\
\texttt{Collection URI} & DOI de la colección en TCIA~\cite{lee2017cbisddsm}. \\
\texttt{License Name} & \textit{Creative Commons Attribution 3.0 Unported}. \\
\bottomrule
\end{tabular}
\end{table}

\subsection{CSV de descripción de casos (a descargar de TCIA)}
Los ficheros \texttt{mass\_/calc\_case\_description\_\{train,test\}\_set.csv}
--- no incluidos en el repositorio; se descargan con \texttt{scripts/download\_cbis\_ddsm.py}
--- sí contienen la etiqueta clínica. La función
\texttt{load\_case\_descriptions} de \texttt{src/data.py} espera, entre otros,
los campos \texttt{patient\_id}, \texttt{breast\_density}
(densidad BI-RADS, 1--4), \texttt{left or right breast}, \texttt{image view}
(CC/MLO), \texttt{abnormality type} (masa o calcificación), \texttt{assessment}
(BI-RADS 0--5), \texttt{subtlety} y, sobre todo, \texttt{pathology}
(\texttt{BENIGN}, \texttt{BENIGN\_WITHOUT\_CALLBACK}, \texttt{MALIGNANT}), que
\texttt{src.data.to\_binary\_label} convierte en la etiqueta binaria 0/1 usada
en todo el trabajo.

\section{Fragmentos de código adicionales}

\subsection{Selección automática de backend de homología persistente}
Como \texttt{giotto-tda} no publica binarios para Python 3.13+
(Sección~\ref{sec:backend-tda}), \texttt{src/topology.py} detecta en tiempo de
importación qué backend usar, sin que el resto del \textit{pipeline} necesite
saberlo:

\begin{lstlisting}[language=Python, caption={Deteccion de backend en src/topology.py (fragmento).}, label={lst:backend-detection}]
try:
    from gtda.homology import CubicalPersistence
    _HAS_GIOTTO = True
except ImportError:
    _HAS_GIOTTO = False

try:
    import gudhi
    _HAS_GUDHI = True
except ImportError:
    _HAS_GUDHI = False

def backend() -> str:
    if _HAS_GIOTTO:
        return "giotto-tda"
    if _HAS_GUDHI:
        return "gudhi"
    return "none"
\end{lstlisting}

\subsection{Imagen de persistencia sin giotto-tda}
Cuando el backend activo es \texttt{gudhi}, la imagen de persistencia se
calcula a mano con un núcleo gaussiano ponderado por persistencia, sobre la
rejilla fija $[0,1]\times[0,1]$ (las imágenes ya están normalizadas por
\texttt{preprocessing.to\_grayscale\_float}, Sección~\ref{cap:metodologia}):

\begin{lstlisting}[language=Python, caption={Imagen de persistencia manual (nucleo gaussiano), src/topology.py (fragmento).}, label={lst:persistence-image-manual}]
grid = np.linspace(0.0, 1.0, n_bins)
birth_grid, pers_grid = np.meshgrid(grid, grid, indexing="ij")
sigma = 1.0 / n_bins

for birth, pers, weight in zip(births, persistences, weights):
    image += weight * np.exp(
        -((birth_grid - birth) ** 2 + (pers_grid - pers) ** 2) / (2 * sigma ** 2)
    )
\end{lstlisting}

Al fijar la rejilla de antemano (en vez de ajustarla, \textit{fit}, por lote
como hace \texttt{giotto-tda}), todas las imágenes del conjunto de datos
producen vectores de la misma longitud y directamente comparables entre sí,
sin necesitar un paso de ajuste previo sobre el conjunto de entrenamiento.

\section{Acceso a los datos longitudinales}
\begin{itemize}
  \item \textbf{EMBED}~\cite{jeong2021embed}: acceso controlado a través de
        Emory; descarga por S3. Cita requerida: \textit{EMory BrEast Imaging
        Dataset (EMBED). HITI Lab, Emory University}.
  \item \textbf{CSAW-CC}~\cite{dataset_csawcc, dembrower2020csaw}: acceso
        controlado a través del Servicio Nacional de Datos de Suecia (SND) y
        la Unidad de Acceso a Datos del Karolinska Institutet; descarga por
        S3. Sus condiciones de uso exigen, en publicaciones, invitar como
        coautores a los curadores del conjunto de datos (criterios ICMJE).
        El 30\,\% de la cohorte no se distribuye y se reserva para
        validaciones en colaboración.
\end{itemize}

\section{Correspondencia entre resultados y \textit{scripts}}
\begin{table}[H]
\centering
\caption{Cada resultado de la memoria y el \textit{script} que lo produce.
Los ficheros de resultados están en \texttt{results/}.}
\label{tab:reproducibilidad}
\small
\begin{tabular}{@{}p{5.6cm}p{8.6cm}@{}}
\toprule
\textbf{Resultado} & \textbf{\textit{Script}} \\
\midrule
Clasificación CBIS-DDSM (Tablas~\ref{tab:resultados}, \ref{tab:refinamiento}) & \texttt{run\_pipeline.py}, \texttt{fuse\_and\_evaluate.py} \\
Interpretabilidad (Tablas~\ref{tab:xai-acuerdo}--\ref{tab:xai-coste}) & \texttt{explainability\_benchmark.py} \\
Tres diseños (Tabla~\ref{tab:tres-disenos}) & \texttt{download\_embed\_longitudinal.py}, \texttt{embed\_longitudinal\_evaluation.py} \\
Resolución frente a dilución & \texttt{embed\_global\_resolution.py} \\
Sitio de la lesión en EMBED (Tabla~\ref{tab:embed-sitio}) & \texttt{embed\_lesion\_site.py} \\
Validación en CSAW-CC (Tabla~\ref{tab:csaw}) & \texttt{download\_csaw.py}, \texttt{csaw\_lesion\_site.py} \\
Registro entre exploraciones & \texttt{registration\_benchmark.py}, \texttt{registered\_lesion\_site.py} \\
Mama contralateral & \texttt{download\_contralateral.py}, \texttt{contralateral\_site.py} \\
Patrón topológico (Figura~\ref{fig:interpretacion}) & \texttt{interpret\_anticipation.py} \\
Densidad frente a topología (Tabla~\ref{tab:densidad}) & \texttt{density\_baseline.py} \\
Localización (Tabla~\ref{tab:localizacion}) & \texttt{csaw\_localization.py} \\
Transferencia entre cohortes (Tabla~\ref{tab:transferencia}) & \texttt{transfer\_model.py} \\
Topología frente a CNN (Tabla~\ref{tab:cnn}) & \texttt{cnn\_vs\_topology.py} \\
Figuras de resultados & \texttt{figuras\_memoria.py} \\
\bottomrule
\end{tabular}
\end{table}
